\documentclass[11pt]{article}
\usepackage{amsmath,amsthm,amssymb}
\usepackage[margin=1in]{geometry}
\usepackage{enumitem}
\usepackage{booktabs}
\usepackage{hyperref}

\newtheorem{theorem}{Theorem}
\newtheorem{lemma}[theorem]{Lemma}
\newtheorem{proposition}[theorem]{Proposition}
\newtheorem{corollary}[theorem]{Corollary}
\theoremstyle{definition}
\newtheorem{definition}[theorem]{Definition}
\newtheorem{remark}[theorem]{Remark}
\newtheorem{example}[theorem]{Example}

\title{General-$r$ dim-additivity for $B_{\ell+1}^{(\lambda)} V_\lambda$\\
   and the structural closed-form SYT formula}
\author{Clio}
\date{13 May 2026 (evening prove session)}

\begin{document}
\maketitle

\begin{abstract}
The May-13 morning paper on multi-corner dimensions conjectured a
closed-form
\[
   \dim B_{\ell+1}^{(\lambda)} V_\lambda \;=\; f^{\lambda^{(\ge 2)}}
\]
in the stable regime. The May-13 afternoon paper
\texttt{2026-05-13-4-3-1n-lower-bound-closed.tex} proved an
\emph{r=2 dim-additivity meta-theorem} structurally for
length-preserving rank $r = 2$, and Remark~9 of that paper sketched
the $r$-general extension. The present paper formalises the
$r$-general version and uses it to give a structural inductive proof
of the closed-form SYT formula under a clean stability hypothesis.

\textbf{Main theorem (general-$r$ dim-additivity).} Let $\lambda$ be
a partition with $r \ge 1$ length-preserving corners
$c_1, \ldots, c_r$ and column-$1$ bottom corner $c_*$. Suppose the
SRD-style reduction
\[
   B_{\ell+1}^{(\lambda)} V_\lambda
   \;=\;
   \mathcal{O}_\lambda \Big[\, \sum_{i=1}^r
   \Phi_i^\lambda\big(B^{(\mu_i)} V_{\mu_i}\big)\Big]
\]
holds, where $\mu_i := \lambda \setminus \{c_i, c_*\}$,
$\Phi_i^\lambda$ is the cell-pair branching injection, and
$\mathcal{O}_\lambda$ is the outer chain. Suppose further that the
corner-pair $T_{n-1}$-$2$-blocks $\{c_i, c_*\}$ are all
non-degenerate (cells in distinct rows and columns; automatic). Then
\[
   \dim B_{\ell+1}^{(\lambda)} V_\lambda
   \;=\; \sum_{i=1}^r \dim B^{(\mu_i)}_{j_0(\mu_i)} V_{\mu_i}.
\]

\textbf{Corollary (closed-form SYT formula).} If $\lambda$ is in
the \emph{fully stable regime} --- meaning that the SRD-style
reduction above holds at every shape encountered in the recursive
descent from $\lambda$ to base case --- then
\[
   \dim B_{\ell+1}^{(\lambda)} V_\lambda \;=\; f^{\lambda^{(\ge 2)}}.
\]

The proof transcribes the r=2 argument verbatim: the unordered
cell pairs $\{c_i, c_*\}$ are pairwise distinct (they share $c_*$
but differ in the other element), so the seminormal SYT supports
of the $\Phi_i^\lambda$-images are pairwise disjoint; each image
lies in $E_{n-1}^+|_{V_\lambda}$; and the outer chain
$\mathcal{O}_\lambda$ acts injectively on $E_{n-1}^+|_{V_\lambda}$
by adjacent-injectivity. The corollary is a straightforward
induction on $|\lambda|$ using the SYT branching rule.

We verify the theorem on r=3 shapes computationally.
\end{abstract}

\section{Setup}\label{sec:setup}

\paragraph{Notation.} As in the May-13 afternoon paper
\texttt{2026-05-13-4-3-1n-lower-bound-closed.tex} (henceforth
\textbf{LB-closed}). $H_q(S_n)$ is the Iwahori--Hecke algebra over
$\mathbb{Q}(q)$ with generators $T_1, \ldots, T_{n-1}$ obeying
$(T_i - q)(T_i + 1) = 0$ and the standard braid relations. $V_\lambda$
is the seminormal cell module for $\lambda \vdash n$ in asymmetric
Murphy normalisation ($b' = 1$). $E_j^\pm$ denote the $q$- and
$(-1)$-eigenspaces of $T_j$ acting on $V_\lambda$ (or any
$H_q(S_n)$-module). For $1 \le k \le n$,
\[
   R'_k \;:=\; (T_{k-1}{+}1)\cdots(T_1{+}1),
   \qquad
   B^{(\lambda)}_{j_0} V_\lambda \;:=\; R'_{j_0}\cdots R'_n V_\lambda
\]
with $j_0 := \ell + 1$, $\ell := \ell(\lambda)$. We always work at
the sharp index $j_0 = \ell + 1$.

\paragraph{Partition data.} Let $\lambda \vdash n$ have length
$\ell = \ell(\lambda)$ and write $\lambda$ as $(\hat\lambda, 1^q)$
with $\hat\lambda$ the part of $\lambda$ with rows of length $\ge 2$
and $q$ the number of trailing $1$-rows. The
\emph{length-preserving corners} of $\lambda$ are the removable
corners $(r, c) \in \lambda$ with $c \ge 2$; equivalently, the
corners of $\hat\lambda$ that sit in column $\ge 2$. List them as
$c_1, \ldots, c_r$. The \emph{column-$1$ bottom} is
$c_* := (\ell, 1)$, which is a removable corner of $\lambda$ exactly
when $\lambda_\ell = 1$, i.e., $q \ge 1$. Throughout this paper we
assume $q \ge 1$ so that $c_*$ is a corner.

\paragraph{Decoration shape.}
$\lambda^{(\ge 2)} := (\lambda_1 - 1, \lambda_2 - 1, \ldots, \lambda_k - 1)$
where $k$ is the number of rows of $\lambda$ with $\lambda_i \ge 2$;
equivalently, $\lambda^{(\ge 2)}$ is $\lambda$ with column~$1$
deleted. There is an explicit bijection between length-preserving
corners of $\lambda$ and corners of $\lambda^{(\ge 2)}$:
$(r, c) \mapsto (r, c-1)$.

\paragraph{Corner-pair branching injections.}
For each length-preserving corner $c_i$, set
\[
   \mu_i \;:=\; \lambda \setminus \{c_i, c_*\}.
\]
Note $|\mu_i| = n - 2$, $\ell(\mu_i) = \ell - 1$. The
\emph{branching injection}
$\Phi_i^\lambda : V_{\mu_i} \hookrightarrow V_\lambda$ is defined
on the seminormal basis by
\[
   \Phi_i^\lambda(v_{T'}^{(\mu_i)})
   \;=\;
   v_{T_A^i} + \tau_i\, v_{T_B^i},
\]
where $T_A^i, T_B^i \in \mathrm{SYT}(\lambda)$ are the two extensions
of $T'$ obtained by placing $\{n-1, n\}$ at $\{c_i, c_*\}$ in the two
orders, and $\tau_i$ is the Hoefsmit off-diagonal coefficient of the
$T_{n-1}$-2-block at the axial distance
$d_i := \mathrm{ax}(c_i) - \mathrm{ax}(c_*)$ between the two cells
(\textbf{LB-closed}, \S 1; cf.~SRD-331 Lemma~3.2). Since $c_i$ is in
column $\ge 2$ and a non-bottom row while $c_*$ is in column $1$
bottom row, $d_i \ne 0$ and the 2-block is non-degenerate.
$\Phi_i^\lambda(v_{T'}^{(\mu_i)})$ is the $+q$-eigenvector of
$T_{n-1}$ on the 2-block $\langle v_{T_A^i}, v_{T_B^i} \rangle$.

\paragraph{Outer chain.}
\[
   \mathcal{O}_\lambda \;:=\;
   (T_\ell{+}1)(T_{\ell+1}{+}1)\cdots(T_{n-2}{+}1),
\]
an $(n - \ell - 1)$-factor product of $(T_j+1)$'s with $j$ running
from $\ell$ up to $n - 2$.

\paragraph{SRD-style reduction (hypothesis).}
We \emph{assume} for the purposes of Theorem~\ref{thm:r-additivity}
that
\begin{equation}\label{eq:reduction}
   B_{\ell+1}^{(\lambda)} V_\lambda
   \;=\;
   \mathcal{O}_\lambda \Big[\,
      \sum_{i=1}^r \Phi_i^\lambda\big(B^{(\mu_i)} V_{\mu_i}\big)
   \Big].
\end{equation}
This is the form taken on by the SRD reduction when all
``other'' pieces vanish: same-row pieces (contributing the
$+q$-eigenspaces of $T_j$ with $j$ inside row blocks of length $\ge 2$)
and pure-length-preserving pieces (contributing $\Phi_{ij}^\lambda$ at
corner pairs $\{c_i, c_j\}$). Both vanishings have been proved in
specific shape families:
\begin{itemize}[leftmargin=2em,topsep=0.2em,itemsep=0.1em]
\item Same-row dies: proved for $(3, 3, 1^q)$ ($q \ge 2$)
   \texttt{2026-05-13-same-row-dies-331.tex}; for $(4, 2, 1^q)$
   ($q \ge 2$) \texttt{2026-05-13-same-row-dies-421.tex}; and
   generally in the $r = 1$ inductive reduction
   \texttt{2026-05-13-r1-inductive-reduction.tex}.
\item Pure-length-preserving dies: proved for $(4, 3, 1^q)$ via
   $j$-formula leakage on the $\{c_1, c_2\}$-block
   \texttt{2026-05-13-4-3-1n-upper-bound.tex} Lemma~3.5.
\end{itemize}
We treat~\eqref{eq:reduction} as a structural hypothesis on the shape
$\lambda$. The present paper's content is the implication
$(\text{\eqref{eq:reduction}}) \Rightarrow \dim B^{(\lambda)} = \sum_i \dim B^{(\mu_i)}$.

\section{Disjoint cell-pair supports}\label{sec:disjoint}

The key combinatorial lemma generalises \textbf{LB-closed} Lemma~1
from $r = 2$ to arbitrary~$r$.

\begin{lemma}[Disjoint SYT supports]\label{lem:disjoint}
For each $i \in \{1, \ldots, r\}$ and each $v \in V_{\mu_i}$, every
seminormal-basis SYT $T$ in the support of $\Phi_i^\lambda(v) \in V_\lambda$
satisfies $\{T^{-1}(n-1), T^{-1}(n)\} = \{c_i, c_*\}$ as an
unordered pair of cells. Consequently, for distinct $i \ne j$,
\[
   \Phi_i^\lambda(V_{\mu_i}) \;\cap\; \Phi_j^\lambda(V_{\mu_j}) \;=\; \{0\}
\]
inside $V_\lambda$, and the sum
$\sum_{i=1}^r \Phi_i^\lambda(V_{\mu_i})$ is a direct sum.
\end{lemma}

\begin{proof}
By construction (\S\ref{sec:setup}), for a seminormal basis vector
$v_{T'}^{(\mu_i)}$ of $V_{\mu_i}$,
\[
   \Phi_i^\lambda(v_{T'}^{(\mu_i)})
   \;=\;
   v_{T_A^i} + \tau_i\, v_{T_B^i},
\]
where $T_A^i, T_B^i$ are the two SYTs of $\lambda$ that restrict to
$T'$ on cells of $\mu_i$ and place $\{n-1, n\}$ at $\{c_i, c_*\}$ in
the two orders. So both supporting SYTs have $\{T^{-1}(n-1), T^{-1}(n)\}
= \{c_i, c_*\}$. Extending by linearity over $v = \sum_{T'} a_{T'}
v_{T'}^{(\mu_i)} \in V_{\mu_i}$, every SYT in the support of
$\Phi_i^\lambda(v)$ has $\{n-1, n\}$ at $\{c_i, c_*\}$.

For $i \ne j$, the unordered cell pairs $\{c_i, c_*\}$ and
$\{c_j, c_*\}$ are distinct as subsets of cells of $\lambda$: they
share $c_*$ but differ in the other element ($c_i \ne c_j$, since
$i \ne j$ and $c_1, \ldots, c_r$ are distinct length-preserving
corners). Hence no single SYT~$T$ of $\lambda$ has its
$\{T^{-1}(n-1), T^{-1}(n)\}$ pair equal to \emph{both} $\{c_i, c_*\}$
and $\{c_j, c_*\}$ simultaneously; equivalently, the SYT supports of
$\Phi_i^\lambda(V_{\mu_i})$ and $\Phi_j^\lambda(V_{\mu_j})$ are
disjoint as subsets of $\mathrm{SYT}(\lambda)$.

The seminormal basis $\{v_T : T \in \mathrm{SYT}(\lambda)\}$ is
linearly independent. So if $v_i \in \Phi_i^\lambda(V_{\mu_i})$ for
$i = 1, \ldots, r$ and $\sum_{i=1}^r v_i = 0$, then expanding each
$v_i$ in the seminormal basis and using pairwise disjoint supports
forces each $v_i = 0$ separately. Hence
$\sum_i \Phi_i^\lambda(V_{\mu_i}) = \bigoplus_i \Phi_i^\lambda(V_{\mu_i})$.
\end{proof}

\begin{lemma}[$\Phi_i^\lambda$ is injective]\label{lem:phi-inj}
Each $\Phi_i^\lambda : V_{\mu_i} \to V_\lambda$ is injective.
\end{lemma}

\begin{proof}
For distinct seminormal basis vectors $v_{T_1'}^{(\mu_i)}$ and
$v_{T_2'}^{(\mu_i)}$ of $V_{\mu_i}$, their images
$\Phi_i^\lambda(v_{T_1'}^{(\mu_i)}) = v_{T_A^{i,1}} + \tau_i\, v_{T_B^{i,1}}$
and
$\Phi_i^\lambda(v_{T_2'}^{(\mu_i)}) = v_{T_A^{i,2}} + \tau_i\, v_{T_B^{i,2}}$
have disjoint $2$-element supports: $T_A^{i,1}, T_B^{i,1}$ are
extensions of $T_1'$ and $T_A^{i,2}, T_B^{i,2}$ are extensions of
$T_2'$; if $T_A^{i,1} = T_A^{i,2}$ then restricting to cells of
$\mu_i$ gives $T_1' = T_2'$, contradicting $T_1' \ne T_2'$.

Hence $\Phi_i^\lambda$ sends a basis of $V_{\mu_i}$ to vectors with
pairwise disjoint, non-zero seminormal supports, which form a
linearly independent set in $V_\lambda$. So $\Phi_i^\lambda$ is
injective.
\end{proof}

\begin{corollary}[Dimension of the inner direct sum]
\label{cor:inner-dim}
\[
   \dim \Big( \bigoplus_{i=1}^r \Phi_i^\lambda(B^{(\mu_i)} V_{\mu_i}) \Big)
   \;=\; \sum_{i=1}^r \dim B^{(\mu_i)} V_{\mu_i}.
\]
\end{corollary}

\begin{proof}
By Lemma~\ref{lem:disjoint}, the sum is direct, so the dimension
equals $\sum_i \dim \Phi_i^\lambda(B^{(\mu_i)})$. By
Lemma~\ref{lem:phi-inj}, each $\Phi_i^\lambda$ is injective, hence
preserves dimensions: $\dim \Phi_i^\lambda(B^{(\mu_i)}) = \dim B^{(\mu_i)}$.
\end{proof}

\section{Image in $E_{n-1}^+$ and adjacent-injectivity}\label{sec:eplus-adj}

The next two ingredients are unchanged from \textbf{LB-closed}; we
restate them for self-containment.

\begin{lemma}[$\Phi_i^\lambda$ lands in $E_{n-1}^+$;
\textbf{LB-closed} Lemma~2]\label{lem:image-Eplus}
For each $i \in \{1, \ldots, r\}$ and every $v \in V_{\mu_i}$,
$\Phi_i^\lambda(v) \in E_{n-1}^+|_{V_\lambda}$.
\end{lemma}

\begin{proof}
By construction, $\Phi_i^\lambda(v_{T'}^{(\mu_i)})$ is the
$+q$-eigenvector of $T_{n-1}$ on the 2-block $\langle v_{T_A^i},
v_{T_B^i}\rangle$ of $V_\lambda$. The 2-block is non-degenerate
because $c_i$ and $c_*$ lie in distinct rows and distinct columns
($c_i$ is in column $\ge 2$ non-bottom row; $c_*$ is in column $1$
bottom row). Hence the $+q$-eigenvector exists and is unique up to
scalar, and $\Phi_i^\lambda(v_{T'}^{(\mu_i)}) \in E_{n-1}^+$.
Extending by linearity, $\Phi_i^\lambda(V_{\mu_i}) \subseteq E_{n-1}^+$.
\end{proof}

\begin{theorem}[Adjacent injectivity chain;
\textbf{LB-closed} Theorem~4 / dim2-331 Theorem~3.4]\label{thm:adj-inj}
For any $H_q(S_n)$-module $V$ and $a \le b \le n-2$, the operator
\[
   (T_a{+}1)(T_{a+1}{+}1)\cdots(T_b{+}1) \big|_{E_{b+1}^+(V)}
   \;:\;
   E_{b+1}^+(V) \;\longrightarrow\; E_a^+(V)
\]
is injective. In particular, the outer chain
$\mathcal{O}_\lambda = (T_\ell{+}1)\cdots(T_{n-2}{+}1)$ acts
injectively on $E_{n-1}^+|_{V_\lambda}$, with image in
$E_\ell^+|_{V_\lambda}$.
\end{theorem}

\section{Main theorem: general-$r$ dim-additivity}\label{sec:main}

\begin{theorem}[Main: general-$r$ dim-additivity]
\label{thm:r-additivity}
Let $\lambda \vdash n$ have length $\ell$, $r \ge 1$
length-preserving corners $c_1, \ldots, c_r$, and column-$1$ bottom
corner $c_*$ (we assume $\lambda_\ell = 1$). Assume the SRD-style
reduction~\eqref{eq:reduction} holds. Then
\begin{equation}\label{eq:r-additivity}
   \dim B_{\ell+1}^{(\lambda)} V_\lambda
   \;=\;
   \sum_{i=1}^r \dim B^{(\mu_i)}_{j_0(\mu_i)} V_{\mu_i}
\end{equation}
where $\mu_i := \lambda \setminus \{c_i, c_*\}$ and
$j_0(\mu_i) := \ell(\mu_i) + 1 = \ell$.
\end{theorem}

\begin{proof}
Define
\[
   W \;:=\; \sum_{i=1}^r \Phi_i^\lambda\big(B^{(\mu_i)} V_{\mu_i}\big)
   \;\subseteq\; V_\lambda.
\]
We compute $\dim W$ and $\dim \mathcal{O}_\lambda W$.

\textbf{Step 1.} $W$ is a direct sum of dimension
$\sum_{i=1}^r \dim B^{(\mu_i)}$.

By Lemma~\ref{lem:disjoint} applied to the inclusion
$B^{(\mu_i)} V_{\mu_i} \subseteq V_{\mu_i}$, the pieces
$\Phi_i^\lambda(B^{(\mu_i)} V_{\mu_i})$ have pairwise disjoint
seminormal-SYT supports inside $V_\lambda$. Hence they form a direct
sum, and by Corollary~\ref{cor:inner-dim},
$\dim W = \sum_i \dim B^{(\mu_i)}$.

\textbf{Step 2.} $W \subseteq E_{n-1}^+|_{V_\lambda}$.

Each summand $\Phi_i^\lambda(B^{(\mu_i)} V_{\mu_i}) \subseteq
\Phi_i^\lambda(V_{\mu_i}) \subseteq E_{n-1}^+|_{V_\lambda}$ by
Lemma~\ref{lem:image-Eplus}. The sum (direct or otherwise) of
subspaces of $E_{n-1}^+|_{V_\lambda}$ remains in
$E_{n-1}^+|_{V_\lambda}$.

\textbf{Step 3.} $\mathcal{O}_\lambda|_W$ is injective.

By Theorem~\ref{thm:adj-inj}, $\mathcal{O}_\lambda$ is injective
on all of $E_{n-1}^+|_{V_\lambda}$, hence on the subspace $W$.

\textbf{Step 4.} Conclude.

By Step~3, $\dim \mathcal{O}_\lambda W = \dim W$. By the SRD-style
reduction~\eqref{eq:reduction}, $B^{(\lambda)} V_\lambda =
\mathcal{O}_\lambda W$. So
\[
   \dim B^{(\lambda)} V_\lambda \;=\; \dim \mathcal{O}_\lambda W
   \;=\; \dim W \;=\; \sum_{i=1}^r \dim B^{(\mu_i)} V_{\mu_i}.
\]
\end{proof}

\begin{remark}[Reading the proof]\label{rem:reading}
The proof is a strict generalisation of \textbf{LB-closed} Theorem~6
(the $r = 2$ dim-additivity): replace ``$2 + 3$'' by
``$\sum_i \dim B^{(\mu_i)}$'' and ``two corner pairs $\{c_1, c_*\},
\{c_2, c_*\}$'' by ``$r$ corner pairs $\{c_1, c_*\}, \ldots, \{c_r, c_*\}$''
throughout. The combinatorial heart (disjoint unordered cell pairs)
generalises trivially because, while the cell pairs share $c_*$, they
differ in the other element, and the other elements $c_1, \ldots, c_r$
are pairwise distinct.
\end{remark}

\section{Closed-form SYT corollary by induction}\label{sec:syt-corollary}

\begin{definition}[Fully stable regime]\label{def:fully-stable}
A partition $\lambda$ is in the \emph{fully stable regime} if the
SRD-style reduction~\eqref{eq:reduction} holds at $\lambda$ and at
every shape $\lambda'$ encountered in the recursive descent
\[
   \lambda \;\longrightarrow\; \{\mu_i : 1 \le i \le r(\lambda)\}
   \;\longrightarrow\; \{\mu_i' : \ldots\} \;\longrightarrow\;
   \cdots \;\longrightarrow\; \text{(column shape $(1^k)$)}.
\]
Equivalently, at every step, the same-row pieces and the pure
length-preserving $\Phi_{ij}$-pieces vanish.
\end{definition}

\begin{lemma}[Decoration shape under corner-pair removal]\label{lem:decoration-removal}
For $\lambda$ with $\lambda_\ell = 1$ and length-preserving corner
$c_i = (\rho_i, \kappa_i)$ ($\kappa_i \ge 2$), define
$\mu_i := \lambda \setminus \{c_i, c_*\}$. Then
\[
   \mu_i^{(\ge 2)} \;=\; \lambda^{(\ge 2)} \setminus c_i^{(\ge 2)},
\]
where $c_i^{(\ge 2)} := (\rho_i, \kappa_i - 1)$ is the cell of
$\lambda^{(\ge 2)}$ corresponding to $c_i$. Moreover,
$c_i^{(\ge 2)}$ is a removable corner of $\lambda^{(\ge 2)}$, and the
map $c_i \mapsto c_i^{(\ge 2)}$ is a bijection between
length-preserving corners of $\lambda$ and removable corners of
$\lambda^{(\ge 2)}$.
\end{lemma}

\begin{proof}
First, the bijection between length-preserving corners of $\lambda$
and corners of $\lambda^{(\ge 2)}$. A length-preserving corner
$c_i = (\rho_i, \kappa_i) \in \lambda$ has $\kappa_i \ge 2$ and is
removable, i.e., $\lambda_{\rho_i} = \kappa_i$ and $\lambda_{\rho_i + 1}
< \kappa_i$. Then $(\rho_i, \kappa_i - 1)$ is a cell of
$\lambda^{(\ge 2)}$ (since $\kappa_i - 1 \ge 1$), and is removable
there: $(\lambda^{(\ge 2)})_{\rho_i} = \lambda_{\rho_i} - 1 = \kappa_i - 1$,
and $(\lambda^{(\ge 2)})_{\rho_i + 1} = \max(\lambda_{\rho_i + 1} - 1, 0)
< \kappa_i - 1$.

Conversely, a removable corner $(\rho, \kappa) \in \lambda^{(\ge 2)}$
has $(\lambda^{(\ge 2)})_\rho = \kappa$, $(\lambda^{(\ge 2)})_{\rho+1}
< \kappa$. Then $\lambda_\rho = \kappa + 1 \ge 2$, $\lambda_{\rho+1}
\le \kappa$. So $(\rho, \kappa + 1) \in \lambda$ is a removable corner
in column $\ge 2$, i.e., length-preserving.

Now compute $\mu_i^{(\ge 2)}$. $\mu_i$ is $\lambda$ with $c_i$ and
$c_*$ removed:
\[
   (\mu_i)_j = \begin{cases}
      \lambda_j - 1 & j = \rho_i, \\
      \lambda_j & j \ne \rho_i, j \ne \ell, \\
      \lambda_\ell - 1 = 0 & j = \ell.
   \end{cases}
\]
So $\mu_i$ has $\ell - 1$ rows (the last row vanishes), and row
$\rho_i$ has $\kappa_i - 1$ cells.

The decoration shape $\mu_i^{(\ge 2)}$ removes column~$1$ from
$\mu_i$: row $j$ of $\mu_i^{(\ge 2)}$ has $(\mu_i)_j - 1$ cells
when $(\mu_i)_j \ge 2$; rows with $(\mu_i)_j \le 1$ become empty.

For $j = \rho_i$: $(\mu_i^{(\ge 2)})_{\rho_i} = \kappa_i - 2$ (if
$\kappa_i \ge 3$) or 0 (if $\kappa_i = 2$, leaving an empty row).
For $j \ne \rho_i, \ell$: $(\mu_i^{(\ge 2)})_j = \lambda_j - 1
= (\lambda^{(\ge 2)})_j$. For $j = \ell$: $(\mu_i^{(\ge 2)})_\ell = 0$
trivially.

On the other hand, $\lambda^{(\ge 2)} \setminus c_i^{(\ge 2)}$ removes
the cell $(\rho_i, \kappa_i - 1)$ from $\lambda^{(\ge 2)}$. Row
$\rho_i$ becomes $\kappa_i - 2$ cells (if $\kappa_i \ge 3$) or
becomes empty (if $\kappa_i = 2$). Other rows unchanged.

So $\mu_i^{(\ge 2)} = \lambda^{(\ge 2)} \setminus c_i^{(\ge 2)}$
matches row-by-row.
\end{proof}

\begin{corollary}[Closed-form SYT formula]\label{cor:syt-formula}
If $\lambda$ is in the fully stable regime
(Definition~\ref{def:fully-stable}), then
\[
   \dim B_{\ell+1}^{(\lambda)} V_\lambda \;=\; f^{\lambda^{(\ge 2)}}.
\]
\end{corollary}

\begin{proof}
By strong induction on $n = |\lambda|$.

\emph{Base case.} If $\lambda$ has no length-preserving corners
($r = 0$), then $\lambda$ is a single column $(1^n)$ and
$\lambda^{(\ge 2)} = \emptyset$. Here $j_0 = \ell + 1 = n + 1 > n$,
so $B^{(\lambda)} = R'_{n+1} R'_{n+2} \cdots R'_n$ is the
\emph{empty product} (the index range is empty), equal to the
identity. Hence $\dim B^{(\lambda)} V_\lambda = \dim V_{(1^n)} = 1
= f^{\emptyset}$. \emph{Note.} In this case the SRD-style reduction
is vacuous: there are no $\mu_i$'s.

\emph{Inductive step.} Suppose $\lambda$ has $r \ge 1$ length-pres
corners. By Theorem~\ref{thm:r-additivity} (applicable since the
fully-stable hypothesis gives~\eqref{eq:reduction}),
\[
   \dim B^{(\lambda)} V_\lambda \;=\; \sum_{i=1}^r \dim B^{(\mu_i)} V_{\mu_i}.
\]
Each $\mu_i$ has $|\mu_i| = n - 2 < n$, and the fully-stable
regime of $\lambda$ implies the fully-stable regime of each $\mu_i$
(the descent tree of $\mu_i$ is a subtree of the descent tree of
$\lambda$). By induction hypothesis,
$\dim B^{(\mu_i)} V_{\mu_i} = f^{\mu_i^{(\ge 2)}}$.

By Lemma~\ref{lem:decoration-removal},
$\mu_i^{(\ge 2)} = \lambda^{(\ge 2)} \setminus c_i^{(\ge 2)}$, where
$c_i^{(\ge 2)}$ ranges over all removable corners of
$\lambda^{(\ge 2)}$ as $i$ ranges over $\{1, \ldots, r\}$. By the
SYT branching identity,
\[
   \sum_{c \text{ corner of } \lambda^{(\ge 2)}} f^{\lambda^{(\ge 2)} \setminus c}
   \;=\; f^{\lambda^{(\ge 2)}}.
\]
Combining,
\[
   \dim B^{(\lambda)} V_\lambda \;=\; \sum_{i=1}^r f^{\mu_i^{(\ge 2)}}
   \;=\; \sum_{c \text{ corner of } \lambda^{(\ge 2)}}
   f^{\lambda^{(\ge 2)} \setminus c} \;=\; f^{\lambda^{(\ge 2)}}.
\]
\end{proof}

\begin{remark}[On the fully-stable hypothesis]\label{rem:stability}
The fully-stable regime hypothesis (Def.~\ref{def:fully-stable}) is
the substantive condition under which the closed-form formula holds.
Concretely, it requires:
\begin{itemize}[leftmargin=2em,topsep=0.2em,itemsep=0.1em]
\item At every shape $\lambda'$ in the descent tree, the
   same-row pieces vanish.
\item At every shape $\lambda'$, the pure length-pres
   $\Phi_{ij}^{\lambda'}$ pieces (corresponding to corner pairs
   $\{c_i, c_j\}$ with $c_*$ \emph{not} in the pair) vanish.
\end{itemize}
These vanishings have been proved for various individual shape
families (cf.~the listed papers in \S\ref{sec:setup}). A general
structural proof of stability has not yet been given; the
shape-dependent threshold $q_0(\widehat\lambda)$ documented in the
May-13 \texttt{2026-05-13-decoration-iso-failed.tex} paper indicates
that stability is genuinely shape-sensitive at small $q$.
The corollary therefore reads:
\emph{whenever stability holds, the SYT formula follows}.
\end{remark}

\section{Computational verification on r=3 shapes}\label{sec:verify}

We verify Theorem~\ref{thm:r-additivity} computationally on $r = 3$
shapes via the dim-recursion
\[
   \dim B^{(\lambda)} \;=\;
   \dim B^{(\mu_1)} + \dim B^{(\mu_2)} + \dim B^{(\mu_3)}.
\]

\paragraph{Setup.} Mod-$P$ arithmetic at $P = 100\,003$, $q = 11$
(matches \texttt{compute\_dim\_B.py}).

\paragraph{Test family $(4, 3, 2, 1^q)$.}
$\hat\lambda = (4, 3, 2)$, $\hat\lambda^{(\ge 2)} = (3, 2, 1)$,
$f^{(3,2,1)} = 16$. Three length-preserving corners:
$c_1 = (1, 4)$, $c_2 = (2, 3)$, $c_3 = (3, 2)$. Inner shapes
and inner decoration shapes:
\begin{align*}
   \mu_1 &= (3, 3, 2, 1^{q-1}), & \mu_1^{(\ge 2)} &= (2, 2, 1), & f^{(2,2,1)} &= 5,\\
   \mu_2 &= (4, 2, 2, 1^{q-1}), & \mu_2^{(\ge 2)} &= (3, 1, 1), & f^{(3,1,1)} &= 6,\\
   \mu_3 &= (4, 3, 1^{q}),       & \mu_3^{(\ge 2)} &= (3, 2),   & f^{(3,2)}   &= 5.
\end{align*}
Predicted by Theorem~\ref{thm:r-additivity}: $\dim B^{(\lambda)} = 5 + 6 + 5 = 16$, matching $f^{(3,2,1)}$ via SYT branching.

\paragraph{Computational results.} Mod-$P$ rank at $P = 100\,003$,
$q = 11$ (computational $q$, not to be confused with the trailing
$1$-row count, which we now call \(t\) for clarity in the table):

\begin{center}\renewcommand{\arraystretch}{1.15}
\begin{tabular}{cccccc}
\toprule
$t$ & $n$ & $\dim V_\lambda$ & $\dim B^{(\lambda)}$ & target $16$ & verdict \\
\midrule
$0$ & $9$  & $168$  & $16$ & $16$ & \emph{matches} \\
$1$ & $10$ & $768$  & $24$ & $16$ & non-stable \\
$2$ & $11$ & $2310$ & $40$ & $16$ & non-stable \\
$3$ & $12$ & $5632$ & $\mathbf{16}$ & $16$ & \checkmark{} stable \\
\bottomrule
\end{tabular}
\end{center}

\paragraph{Verification of additivity at $t = 3$.}
Computing each inner $\dim B^{(\mu_i)}$ separately at $t = 3$
(so inner shapes have trailing 1-row count $t - 1 = 2$ for $\mu_1,
\mu_2$ and $t = 3$ for $\mu_3$):
\begin{center}\renewcommand{\arraystretch}{1.15}
\begin{tabular}{cccc}
\toprule
shape & $n$ & $\dim B$ & target ($f^{\mu^{(\ge 2)}}$) \\
\midrule
$\mu_1 = (3, 3, 2, 1^2)$ & $10$ & $5$ & $f^{(2,2,1)} = 5$ \checkmark \\
$\mu_2 = (4, 2, 2, 1^2)$ & $10$ & $6$ & $f^{(3,1,1)} = 6$ \checkmark \\
$\mu_3 = (4, 3, 1^3)$    & $10$ & $5$ & $f^{(3,2)}   = 5$ \checkmark \\
\midrule
sum                       &      & $16$ & matches $\dim B^{(\lambda)} = 16$ \checkmark \\
\bottomrule
\end{tabular}
\end{center}

The stability threshold for this family is $t \ge 3$, sharper than
the crude $t \ge |\lambda^{(\ge 2)}| = 6$ bound.

\paragraph{Curious observation at $t = 0$.}
At $t = 0$ (the no-trailing-1-rows case $\lambda = \hat\lambda = (4, 3, 2)$,
$n = 9$), the computed $\dim B = 16$ \emph{matches} the closed-form
target, even though our framework formally requires $t \ge 1$ for the
column-$1$ corner $c_*$ to exist. The matching-failing-matching
pattern $\{0 : \checkmark, 1 : \times, 2 : \times, 3 : \checkmark, \ldots\}$
echoes the non-monotone-in-$t$ behaviour observed for $\widehat\lambda = (4, 2)$
in \texttt{2026-05-13-decoration-iso-failed.tex}. The $t = 0$ match
suggests an independent reduction (with a different cell structure)
operates at the no-trailing-rows boundary. We leave this for future
work.

Scripts: \texttt{\textasciitilde/projects/scratch/2026-05-13-r-additivity/test\_r3\_fast.py},
\texttt{compute\_dim\_B\_fast.py}.

\section{Honest status}\label{sec:status}

\paragraph{What is proved structurally.}
Theorem~\ref{thm:r-additivity} (general-$r$ dim-additivity) is proved
structurally, modulo only the SRD-style reduction
hypothesis~\eqref{eq:reduction}. The proof transcribes the r=2 proof
of \textbf{LB-closed} verbatim, with the combinatorial heart
(disjoint cell pairs) generalising trivially.

\paragraph{What is conditional.}
Corollary~\ref{cor:syt-formula} (closed-form SYT formula) requires
the fully-stable regime hypothesis at all descent shapes. Stability
is shape-dependent and is itself a substantial structural fact; it
has been proved in individual shape families but not yet in full
generality.

\paragraph{What this resolves.}
The Remark~9 of \textbf{LB-closed} sketched the general-$r$
extension; the present paper formalises it. The closed-form SYT
formula, which was conjectural in
\texttt{2026-05-13-multi-corner-dimension.tex}, becomes a
\emph{structural} theorem (under the fully-stable hypothesis) rather
than a numerical coincidence.

\paragraph{Comparison to the failed decoration iso.}
The naive Hecke-iso upgrade was refuted in
\texttt{2026-05-13-decoration-iso-failed.tex}: there is no parabolic
embedding of $H_q(S_{n-\ell})$ inside $H_q(S_n)$ acting on
$B^{(\lambda)} V_\lambda$ as $V_{\lambda^{(\ge 2)}}$. The present
paper achieves the next best thing: a \emph{numerical} equality of
dimensions, derived structurally via direct decomposition of
$V_\lambda$, without claiming any rep-theoretic content beyond
counting.

\paragraph{Next steps.}
\begin{itemize}[leftmargin=2em,topsep=0.2em,itemsep=0.1em]
\item Prove the SRD-style reduction~\eqref{eq:reduction} for all
   stable shapes structurally (currently only proved in special
   families).
\item Characterise the shape-dependent stability threshold
   $q_0(\widehat\lambda)$ (the failed-iso paper's open problem).
\item Determine whether $B^{(\lambda)} V_\lambda$ carries a Hecke
   action via exotic operators (Jucys--Murphy, intertwiners) that
   upgrade the numerical iso to a rep-theoretic one.
\end{itemize}

\end{document}
